% ======================================================================
%  Front Matter of The Actualization Theory
%  AT-MONO115 · Canonical Monograph (v2.0) · Zenodo-ready
%  Includes: front cover, title page, copyright, abstract, executive
%  summary, preface, table of contents.
%  Assembly only — no new physics, no new derivations.
% ======================================================================

% ---- Title page (also serves as front cover) ------------------------------
\thispagestyle{empty}
\begin{center}
  \vspace*{2.5cm}
  {\Huge\bfseries The Actualization Theory}\\[1em]
  {\LARGE A Reconstruction of Physics from\\ Difference, Actualization and Spectrum}\\[2.5cm]
  {\Large Fabrice Wieser}\\[1mm]
  {\small Independent Researcher and Full Stack Software Engineer}\\[0.5mm]
  {\small \url{https://github.com/MagusDraconis/AT}}\\[2.5cm]
  {\large The Actualization Theory (AT) --- v2.0 (canonical)}\\[1cm]
  {\large \today}\\[1.5cm]
  {\large Zenodo edition}
  \vfill
\end{center}
\clearpage

% ---- Copyright page --------------------------------------------------------
\thispagestyle{empty}
\begin{center}
  \vspace*{4cm}
  {\large The Actualization Theory}\\[0.6em]
  {\normalsize A Reconstruction of Physics from Difference, Actualization and Spectrum}\\[2em]
  {\normalsize Copyright \copyright\ \the\year\ Fabrice Wieser}\\[1.5em]
  {\small All rights reserved.}
  \vfill
  {\small Version 2.0 (canonical) --- \today}\\
  {\small Compiled from the canonical AT-QG research program (QG0--QG321).}
  \vfill
\end{center}
\clearpage

% ---- Abstract ---------------------------------------------------------------
\chapter*{Abstract}
\label{fm:abs:abstract}
\addcontentsline{toc}{chapter}{Abstract}

The Actualization Theory reconstructs physics from the primitive Difference and the tensor
reference $\eta$, through the canonical hierarchy

\[
\text{Difference} \;\longrightarrow\; \text{Actualization}
\;\longrightarrow\; \text{Inevitable Spectrum} \;\longrightarrow\; \text{Physics}.
\]

Difference is the fundamental boundary: the irreducible primitive whose count-conservation
law is its definitional identity, and whose two faces --- the scalar counting measure
$\rho$ and the tensor field $\psi$ --- carry, respectively, the spectral readout of all
physics and the gravitational content of the theory. Actualization is Difference's derived
count-producing process, whose convergence is the closure boundary $N=96$. The spectrum is
the inevitable output of the actualization attractor: the unique eigenspectrum of the
converged network, from which every physics sector is read.

The monograph derives, from this foundation: quantum mechanics ($|\psi|^2=\rho$,
measurement as actualization), gravity and spacetime (emergent geometry, native dynamics),
the standard model (fermion sectors, gauge structure, couplings, mixing, weak scale,
precision predictions $S,T,U$, the electron g-2, the Majorana character), and cosmology
(the cosmological constant, the density fractions, the CMB acoustic structure, structure
formation). The boundaries of the theory are documented explicitly: Difference as the
ontological boundary, the unresolved status of $\psi$, the numerical boundary $\pi$, the
Bekenstein $2\pi$ factor, and the hosted standard-model dynamics. The theory claims no open
derivation frontier; its frontier is the experimental validation of its pre-registered
predictions.

% ---- Executive summary --------------------------------------------------------
\chapter*{Executive Summary}
\label{fm:abs:executive-summary}
\addcontentsline{toc}{chapter}{Executive Summary}

\textbf{The foundation.} The Actualization Theory uses exactly two primitives,
$\{\text{Difference},\eta\}$. Difference is the fundamental boundary --- the irreducible
counting difference from a uniform background --- and $\eta$ is the tensor reference against
which the tensor face of Difference is read. Actualization, the count-producing process, is
derived from Difference, not a third primitive.

\textbf{The spectrum.} The D96 spectrum is the inevitable output of the actualization
attractor: one zero mode and 95 positive modes, with multiplicities
$[42\times2,5,6]$, moments $\Sigma m=95$, $\Sigma\sqrt{m}=64.08$, $\Sigma m^2=229$, the
occupancy moment $\mathrm{occMom}=1900.25$, and the span $6.40$. Every derived constant
traces to this spectrum.

\textbf{The physics.} Quantum mechanics is emergent ($|\psi|^2=\rho$, measurement is
actualization). Gravity is emergent (the coupling from the D96 spectral content, the tensor
sector from $\psi$ against $\eta$, spacetime from the counting measure). The standard model
is a set of derived reads of the D96 moments: the $1+3+8$ gauge structure, the couplings,
the CKM and PMNS matrices, the weak scale, the Higgs, the neutrino masses, and the
precision predictions. Cosmology is a spectrum readout: the cosmological constant is the
residual actualization pressure, the density fractions are the octave-record information,
and the CMB acoustic structure is the spectral harmonics.

\textbf{The boundaries.} The theory documents its boundaries honestly: Difference as the
ontological boundary, the unresolved ontological status of $\psi$, the numerical boundary
$\pi$, the Bekenstein $2\pi$ factor, and the hosted standard-model dynamics. These are
limits of the framework, not derivation failures.

\textbf{The frontier.} The theory claims no open derivation frontier. Its frontier is the
experimental validation of the derived predictions: the oblique parameters $S,T,U$, the
electron g-2 $a_e$, the Majorana character $0\nu\beta\beta$, and the pre-registered
predictions P1 (106 GeV), P2 ($0\nu\beta\beta$), and the P3 sector ladder.

% ---- Preface ------------------------------------------------------------------
\chapter*{Preface}
\label{fm:abs:preface}
\addcontentsline{toc}{chapter}{Preface}

This monograph is the canonical v2.0 record of The Q-Model (AT), now presented under the
name \emph{The Actualization Theory}. It consolidates the end-state of a research program
that began with the question of whether observable physics can be reconstructed from a
smaller primitive foundation, and converged on the canonical hierarchy
Difference $\to$ Actualization $\to$ Spectrum $\to$ Physics.

The monograph is a \emph{foundation monograph}: it presents the primitives, the derivation
hierarchy, the physics, the boundaries, and the validation frontier. It is written to be
referee-safe --- every claim is either a derived result with its source, a documented
boundary, or an experimental-validation item awaiting its experiment. No claim exceeds the
accepted derivations of the canonical program.

The structure follows the physics-focused plan: the foundation (Parts I--II), the spectrum
(Part III), the physics (Part IV), and the boundary layer and frontier (Part V), followed
by the general conclusion and the appendices. The universality research program --- the
operator basis across organized domains, the lock law, and the organization structure ---
is treated as future work in the appendix, outside the core physics derivation.

% ---- Table of contents ---------------------------------------------------------
\tableofcontents
\clearpage
